% =====================================================================
%  tr03a.tex —— 原文 p.8（tr-03）
%  内容：定理 3（Bott 周期性 $\Omega^2BU\simeq BU\times\Z$）、$\Omega$-预谱与 $\widetilde K^n$、
%        §3.1 另一版本的 Bott 周期性（定理 4、定理 5）、§4 $K$ 理论的一个应用（定理 6）
% =====================================================================
% ===== 原文 p.8 =====
% [原文 p.8]

\begin{thm}{3}
  （Bott 周期性，[AGP] 9.5.1 及附录 B）存在同伦等价
  \[ \Omega^2 BU\simeq BU\times\Z . \]
  [分类空间 $BU$ 的同伦群以 2 为周期
  \[
  \pi_{i+2}(BU)=\pi_i(\Omega^2BU)=\pi_i(BU\times\Z)=
  \begin{cases}
    \Z & i=0,\\
    \pi_i(BU)=0 & i\ge 1,
  \end{cases}
  \]
  故称「周期性」。]
\end{thm}
% [Tutor Note: 定理 3 的方括号说明照原文译出，但第二行 $\pi_i(BU)=0$（$i\ge 1$）是原文笔误。1100 dpi 裁图逐字复核：印刷确为 "$\pi_i(BU)=0\quad i\ge 1,$"。正确形式是 $\pi_i(BU)$（不含 $=0$）：由 $\pi_i(BU\times\Z)=\pi_i(BU)\oplus\pi_i(\Z)$ 与 $\pi_i(\Z)=0$（$i\ge 1$）可知，$i\ge 1$ 时它只能是 $\pi_i(BU)$。反向自检：若该行真为 0，则 $\pi_2(BU)=0$；但同页下文说 $\widetilde{K}(S^2)\cong\Z$（由 $H-1$ 生成），而 $\pi_2(BU)=[S^2,BU]=\widetilde{K}(S^2)\cong\Z$，矛盾。正确写法应为 $\pi_{i+2}(BU)=\pi_i(BU)$：$i\ge 1$ 时它就是 $\pi_i(BU)$；$i=0$ 那一行给的是 $\pi_2(BU)=\pi_0(BU\times\Z)=\pi_0(BU)\oplus\Z=\Z$（$BU$ 连通故 $\pi_0(BU)=0$），这正是「周期为 2」的内容。]

Bott 周期性表明，我们有一个由
\[ E_{2n}=BU\times\Z,\quad E_{2n+1}=\Omega BU, \]
构成的 $\Omega$-预谱（$\Omega$-prespectrum），从而得到约化上同调理论
\[ \widetilde E^n(X)=[X,E_n]_*=:\widetilde K^n(X) \]
它在紧 Hausdorff 空间 $X$ 上与我们先前的定义一致。
% [Tutor Note: 这里的 $\{E_n\}_{n\in\Z}$ 接的是 p.7 定义 6 的 $\Omega$-预谱（带弱同伦等价 $\epsilon_n:E_n\to\Omega E_{n+1}$）。由定理 3 可直接读出这些结构映射：$\Omega E_{2n+1}=\Omega^2BU\simeq BU\times\Z=E_{2n}$，而 $\Omega E_{2n+2}=\Omega(BU\times\Z)\simeq\Omega BU=E_{2n+1}$（因为 $\Omega\Z\simeq*$），两族都是同伦等价。等价地说 $E_{n+2}\simeq E_n$，故该谱以 2 为周期；$\widetilde E^n(X)=[X,E_n]_*$ 里的下标 $*$ 表示保持基点的同伦类，与定义 5 的约化上同调理论对接。]

\subsection{另一版本的 Bott 周期性}

给定空间 $X$ 与 $Y$，投影 $p_1,p_2:X\times Y\to X,Y$ 给出一个乘法
\[
K(X)\otimes K(Y)\xrightarrow{\ \mu\ }K(X\times Y),\quad \mu(a\otimes b)=p_1^*(a)p_2^*(b)
\]
或在约化的情形
\[ \widetilde{K}(X)\otimes\widetilde{K}(Y)\xrightarrow{\ \tilde\mu\ }\widetilde{K}(X\wedge Y). \]
取 $Y=S^2$ 便得到同构。

\begin{thm}{4}
  （[VBKT] 2.2，[Bott] 附录 I）通过上面定义的外积 $\mu$、$\tilde\mu$，我们有同构：
  \[
  K(X)\otimes K(S^2)\xrightarrow{\ \cong\ }K(X\times S^2),\quad
  \widetilde{K}(X)\otimes\widetilde{K}(S^2)\xrightarrow{\ \cong\ }\widetilde{K}(X\wedge S^2)\cong\widetilde{K}(\Sigma^2X).
  \]
\end{thm}

知道 $\widetilde{K}(S^2)\cong\Z$（由 $H-1$ 生成）就给出 Bott 周期性的另一个版本。
% [Tutor Note: 定理 4 的两个同构分属未约化情形（$K(X)\otimes K(S^2)\to K(X\times S^2)$）与约化情形（$\widetilde{K}(X)\otimes\widetilde{K}(S^2)\to\widetilde{K}(X\wedge S^2)$），不能互相代入；右端写 $\widetilde{K}(X\wedge S^2)\cong\widetilde{K}(\Sigma^2X)$ 是因为 $\Sigma^2X=X\wedge S^2$。定理 5 把整个 $\widetilde{K}(S^2)$ 换成一个生成元 $H-1$（$H$ 为 §1 例中 $P^1(\C)$ 上的典范丛，见 p.1 与本页上一段的提醒），箭头上的 $\otimes(H-1)$ 就是「乘以 $H-1$」这个映射，与约化外积 $\tilde\mu$ 复合即得每个 $X$ 上周期为 2 的同构。]

\begin{thm}{5}
  （Bott 周期性，[VBKT] 2.11）存在同构
  \[
  \widetilde{K}(X)\xrightarrow{\ \otimes(H-1)\ }\widetilde{K}(X)\otimes\widetilde{K}(S^2)
  \xrightarrow{\ \tilde\mu\ }\widetilde{K}(X\wedge S^2)\cong\widetilde{K}(\Sigma^2X),
  \]
  使得 $X$ 的约化 $K$ 理论在悬挂下是 2-周期的。
\end{thm}

\section{$K$ 理论的一个应用}

尽管此前已用其他技术证明过，但利用 $K$ 理论可以给出下述定理的一个（相对）简单的证明，参见 [VBKT] 第 2.3 节或 [AGP] 第 10 章。

\begin{thm}{6}
  仅有的 $H$-空间球面是
  \[ S^0,\ S^1,\ S^3,\ S^7, \]
  它们来自 $\R$、$\C$、$\mathbb{H}$ 与 $\mathbb{O}$ 中范数为一的元素上的代数结构。
\end{thm}
% [Tutor Note: 定理 6 的末句只说明了「为什么」$S^0,S^1,S^3,S^7$ 是 $H$-空间（$\R$、$\C$、$\mathbb{H}$、$\mathbb{O}$ 中范数为一的元素分别构成这四个球面，其代数乘法给出 $H$-空间结构）；定理的实质内容是「其余球面都不是 $H$-空间」，原文未给证明，只指向 [VBKT] 第 2.3 节与 [AGP] 第 10 章。p.9 开头的提示说明该证明要用 Adams 运算与分裂原理。]
